% GENERATED FILE. Edit canonical lessons, then run npm run generate:books.
% canonical-source-hash: fnv1a64:20262e62f368c3f9
% canonical-lessons: FR-C210-atteindre, FR-C210-realiser, FR-C210-offrir-un-cadeau, FR-C210-rendre, FR-C210-reconnaitre

\chapter{To reach, To achieve, To give a present, To give back, To recognise}
\label{ch:fr-to-reach-to-achieve-to-give-a-present-to-give-back-to-recognise}
\hlchaptermodality{210}

\begin{chapteropening}
\textbf{By the end of this chapter:} \emph{I can say to reach, to achieve, to give a present, to give back and to recognise.}

Complete the last lesson of chapter 210: I can say to reach, to achieve, to give a present, to give back and to recognise.
\end{chapteropening}

\section[atteindre]{atteindre — to reach}

\label{lesson:FR-C210-atteindre}



\begin{quote}
Before the new one: say the French for to inform, then the French for to complain.
\end{quote}

\subsection*{You'll want to know: atteindre}
\textbf{atteindre} — \textquotedblleft{}to reach\textquotedblright{}.

\begin{cousinweb}
\textbf{atteindre} $\leftarrow$ Old French \emph{ateindre} (EtymDB 2.1) $\leftarrow$ Latin \emph{attingere}, \textquotedblleft{}to touch upon, to reach\textquotedblright{} (Lewis \& Short) — inherited.
\end{cousinweb}

\begin{grammarlens}[title={What you've built}]
One more word for this chapter.
\end{grammarlens}

\subsection*{Your turn}
\begin{itemize}
\raggedright
  \item \emph{Say it:} \emph{atteindre}
  \item \emph{Say it:} \emph{atteindre}, once more
  \item \emph{Say it:} say \emph{atteindre}
  \item \emph{Recall:} say the French for to inform, then the French for to complain, then say \emph{atteindre} again
\end{itemize}

\begin{tcolorbox}[breakable,colback=teal!4,colframe=teal!35!black,title={Before you move on}]
What does atteindre mean? (\textquotedblleft{}To reach\textquotedblright{}.) Say it once more.
\end{tcolorbox}
\section[réaliser]{réaliser — to achieve}

\label{lesson:FR-C210-realiser}



\begin{quote}
Before the new one: say the French for to complain, then the French for to reach.
\end{quote}

\subsection*{You'll want to know: réaliser}
\textbf{réaliser} — \textquotedblleft{}to achieve\textquotedblright{}.

\begin{cousinweb}
\textbf{réaliser} — built from the word for \textquotedblleft{}real\textquotedblright{}, $\leftarrow$ Late Latin \emph{realis}, from \emph{res}, \textquotedblleft{}a thing\textquotedblright{}, with \emph{-iser}: to make real. (EtymDB 2.1)
\end{cousinweb}

\begin{grammarlens}[title={What you've built}]
One more word for this chapter.
\end{grammarlens}

\subsection*{Your turn}
\begin{itemize}
\raggedright
  \item \emph{Say it:} \emph{réaliser}
  \item \emph{Say it:} \emph{réaliser}, once more
  \item \emph{Say it:} say \emph{réaliser}
  \item \emph{Recall:} say the French for to complain, then the French for to reach, then say \emph{réaliser} again
\end{itemize}

\begin{tcolorbox}[breakable,colback=teal!4,colframe=teal!35!black,title={Before you move on}]
What does réaliser mean? (\textquotedblleft{}To achieve\textquotedblright{}.) Say it once more.
\end{tcolorbox}
\section[offrir un cadeau]{offrir un cadeau — to give a present}

\label{lesson:FR-C210-offrir-un-cadeau}



\begin{quote}
Before the new one: say the French for to reach, then the French for to achieve.
\end{quote}

\subsection*{You'll want to know: offrir un cadeau}
\textbf{offrir un cadeau} — \textquotedblleft{}to give a present\textquotedblright{}.

\begin{cousinweb}
\textbf{offrir un cadeau} — \emph{offrir}, \textquotedblleft{}to offer\textquotedblright{}, $\leftarrow$ Latin \emph{offerre} (EtymDB 2.1), \emph{un}, \textquotedblleft{}a\textquotedblright{}, and \emph{cadeau}, \textquotedblleft{}a present\textquotedblright{}.
\end{cousinweb}

\begin{grammarlens}[title={What you've built}]
One more word for this chapter.
\end{grammarlens}

\subsection*{Your turn}
\begin{itemize}
\raggedright
  \item \emph{Say it:} \emph{offrir un cadeau}
  \item \emph{Say it:} \emph{offrir un cadeau}, once more
  \item \emph{Say it:} say \emph{offrir un cadeau}
  \item \emph{Recall:} say the French for to reach, then the French for to achieve, then say \emph{offrir un cadeau} again
\end{itemize}

\begin{tcolorbox}[breakable,colback=teal!4,colframe=teal!35!black,title={Before you move on}]
What does offrir un cadeau mean? (\textquotedblleft{}To give a present\textquotedblright{}.) Say it once more.
\end{tcolorbox}
\section[rendre]{rendre — to give back}

\label{lesson:FR-C210-rendre}



\begin{quote}
Before the new one: say the French for to achieve, then the French for to give a present.
\end{quote}

\subsection*{You'll want to know: rendre}
\textbf{rendre} — \textquotedblleft{}to give back\textquotedblright{}.

\begin{cousinweb}
\textbf{rendre} $\leftarrow$ Vulgar Latin \emph{*rendere} $\leftarrow$ Latin \emph{reddere}, \textquotedblleft{}to give back\textquotedblright{} (EtymDB 2.1) — inherited, its \emph{-n-} taken from the Latin verb for \textquotedblleft{}to take\textquotedblright{}.
\end{cousinweb}

\begin{grammarlens}[title={What you've built}]
One more word for this chapter.
\end{grammarlens}

\subsection*{Your turn}
\begin{itemize}
\raggedright
  \item \emph{Say it:} \emph{rendre}
  \item \emph{Say it:} \emph{rendre}, once more
  \item \emph{Say it:} say \emph{rendre}
  \item \emph{Recall:} say the French for to achieve, then the French for to give a present, then say \emph{rendre} again
\end{itemize}

\begin{tcolorbox}[breakable,colback=teal!4,colframe=teal!35!black,title={Before you move on}]
What does rendre mean? (\textquotedblleft{}To give back\textquotedblright{}.) Say it once more.
\end{tcolorbox}
\section[reconnaître]{reconnaître — to recognise}

\label{lesson:FR-C210-reconnaitre}



\begin{quote}
Before the new one: say the French for to give a present, then the French for to give back.
\end{quote}

\subsection*{You'll want to know: reconnaître}
\textbf{reconnaître} — \textquotedblleft{}to recognise\textquotedblright{}.

\begin{cousinweb}
\textbf{reconnaître} $\leftarrow$ Old French \emph{reconoistre} $\leftarrow$ Latin \emph{recognoscere}, \textquotedblleft{}to know anew\textquotedblright{} (EtymDB 2.1) — inherited: \emph{connaître} with \emph{re-}.
\end{cousinweb}

\begin{grammarlens}[title={What you've built}]
One more word for this chapter.
\end{grammarlens}

\subsection*{Your turn}
\begin{itemize}
\raggedright
  \item \emph{Say it:} \emph{reconnaître}
  \item \emph{Say it:} \emph{reconnaître}, once more
  \item \emph{Say it:} say \emph{reconnaître}
  \item \emph{Recall:} say the French for to give a present, then the French for to give back, then say \emph{reconnaître} again
\end{itemize}

\begin{tcolorbox}[breakable,colback=teal!4,colframe=teal!35!black,title={Before you move on}]
What does reconnaître mean? (\textquotedblleft{}To recognise\textquotedblright{}.) Say it once more.
\end{tcolorbox}
